% Generated by experiments/paper_materials.py from authoritative formal outputs.
\begin{table}[t]
\centering
\scriptsize
\setlength{\tabcolsep}{3pt}
\renewcommand{\arraystretch}{1.12}
\resizebox{\columnwidth}{!}{%
\begin{tabular}{@{}cccccccc@{}}
\hline
Seed & $\|R\|_2$ & $\|\nabla_cJ\|_2$ & $\|\nabla_aJ\|_2$ & $\Delta_c$ & $\Delta_a$ & $\lambda_{\min}(H_c)$ & $\lambda_{\min}(H_a)$ \\
\hline
0 & 4.365 & $5.28\times10^{-15}$ & $6.93\times10^{-14}$ & 0.8707 & -0.0305 & 0.8856 & 0.0002 \\
1 & 4.222 & $4.89\times10^{-13}$ & $7.56\times10^{-13}$ & 0.9316 & 0.0011 & 0.9514 & 0.0062 \\
2 & 7.183 & $3.59\times10^{-11}$ & $9.29\times10^{-11}$ & 0.5462 & -0.9645 & 0.7298 & 0.0098 \\
3 & 5.371 & $1.94\times10^{-15}$ & $6.16\times10^{-15}$ & 0.8514 & -0.3651 & 0.9073 & 0.0063 \\
4 & 6.810 & $1.30\times10^{-12}$ & $3.44\times10^{-12}$ & 0.5519 & -2.4977 & 0.7800 & 0.0002 \\
\hline
\end{tabular}
}
\caption{Coordinate dependence of the sufficient curvature inequality at the same physical endpoints.  Subscript $c$ denotes the fixed QR coordinates used for optimization and $a$ the raw feature coefficients; $\Delta$ is the sufficient curvature margin.  The QR margin is positive for all five seeds, whereas the raw-coordinate margin is positive for one.  Both coordinate Hessians remain positive definite in every row, so a negative margin means failure of this sufficient test, not a saddle point.}
\label{tab:coordinate-certificate}
\end{table}
